\subsection{Markov chains, invariant distributions, and ergodicity}
\label{subsec:markov-chains-kernels}

In many problems our goal is to compute expectations under a distribution $\pi$ on $\R^d$,
\[
\mu \coloneqq \E_{X\sim \pi}[g(X)],
\]
where $\pi$ may be known only up to a normalizing constant.
If we can draw i.i.d.\ samples from $\pi$, then the basic Monte Carlo estimator is immediate.
However, in moderate or high dimensions, designing good i.i.d.\ proposals can be difficult; for instance, rejection sampling can have vanishing acceptance probabilities and importance sampling can suffer from weight degeneracy (see \Cref{subsec:mc-generating}).

Markov chain Monte Carlo (MCMC) takes a different approach.
Instead of sampling i.i.d., we simulate a Markov chain whose \emph{marginal} distribution converges to $\pi$.
The resulting samples are correlated, but if the chain mixes well then empirical averages of $g(X_n)$ can still be accurate.
Thus an MCMC method is a choice of a one-step transition rule $K$ that is (i) easy to simulate, (ii) has invariant distribution $\pi$, and (iii) moves through the state space efficiently.

\begin{definition}[Transition kernel]
Let $\mathcal{X}=\R^d$ and let $\mathcal{B}(\mathcal{X})$ denote its Borel $\sigma$-algebra.
A \emph{transition kernel} is a function
\[
K:\mathcal{X}\times\mathcal{B}(\mathcal{X})\to[0,1],\qquad (x,A)\mapsto K(x,A),
\]
such that:
(i) for each fixed $x\in\mathcal{X}$, the map $A\mapsto K(x,A)$ is a probability measure on $\mathcal{B}(\mathcal{X})$, and
(ii) for each fixed $A\in\mathcal{B}(\mathcal{X})$, the map $x\mapsto K(x,A)$ is measurable.
\end{definition}

The interpretation is that $K(x,A)$ is the probability that one step of the chain moves into the set $A$ when the current state is $x$.
We often write $K(x,dy)$ as a notational reminder that, for each fixed $x$, $K(x,\cdot)$ is a measure on $\mathcal{X}$.

If $K(x,\cdot)$ is absolutely continuous with respect to Lebesgue measure, we write $k(y\mid x)$ for a conditional density and
\[
K(x,A) = \int_A k(y\mid x)\,dy.
\]
More generally, for any measurable $f:\mathcal{X}\to\R$ with $\int |f(y)|\,K(x,dy)<\infty$, define the associated Markov operator
\[
(Kf)(x) \coloneqq \int_{\mathcal{X}} f(y)\,K(x,dy).
\]
Then $(Kf)(x)$ is the conditional expectation $\E[f(X_{n+1})\mid X_n=x]$ when $(X_n)$ has transition kernel $K$.

\begin{definition}[Markov chain with kernel $K$]
A sequence $(X_n)_{n\ge 0}$ is a (time-homogeneous) Markov chain with transition kernel $K$ if for every $n\ge 0$ and every $A\in\mathcal{B}(\mathcal{X})$,
\[
\mathbb{P}(X_{n+1}\in A\mid X_0,\dots,X_n)=\mathbb{P}(X_{n+1}\in A\mid X_n)=K(X_n,A)
\]
almost surely.
\end{definition}

Equivalently, conditional on $X_n=x$, we draw the next state as
\[
X_{n+1}\mid X_n=x \sim K(x,\cdot).
\]
An initial distribution for $X_0$, together with $K$, determines the joint law of $(X_0,X_1,\dots)$.
We focus on discrete-time, time-homogeneous chains because they are directly implementable, though many MCMC methods are motivated by continuous-time Markov processes and their discretizations.

\subsubsection{Transition kernels and iterated kernels}
\label{subsubsec:iterated-kernels}
Let $K$ be a transition kernel.
Define the $m$-step kernel $K^m$ recursively by $K^1=K$ and
\[
K^{m+1}(x,A) \coloneqq \int_{\mathcal{X}} K^m(y,A)\,K(x,dy).
\]
Then $K^m(x,A)=\mathbb{P}(X_{n+m}\in A\mid X_n=x)$, and the recursion is the discrete-time Chapman--Kolmogorov equation.
If $\mu$ is a distribution for $X_0$, the marginal distribution after $m$ steps is the measure
\[
(\mu K^m)(A)\coloneqq \int_{\mathcal{X}} K^m(x,A)\,\mu(dx).
\]
Thinking in these terms makes the MCMC objective precise: we want $\mu K^m$ to approach $\pi$ as $m$ grows, even when $\mu$ is far from $\pi$.
Continuous-time Markov processes use the same notation with a time index:
for a time-homogeneous process \((X_t)_{t\ge 0}\), one writes
\[
  K_t(x,A)\coloneqq \mathbb{P}(X_{s+t}\in A\mid X_s=x),
  \qquad
  t\ge 0,
\]
which does not depend on \(s\).
Then \(K_0(x,A)=\1\{x\in A\}\), and the Chapman--Kolmogorov equation becomes
\[
  K_{s+t}(x,A)=\int_{\mathcal{X}} K_t(y,A)\,K_s(x,dy),
  \qquad
  s,t\ge 0.
\]
This is the continuous-time analogue of the discrete iterated-kernel recursion above.

\begin{example}[Gaussian AR(1) chain]
\label{ex:ar1-kernel}
Fix $\rho\in(-1,1)$ and let $(\varepsilon_n)_{n\ge 1}$ be i.i.d.\ $\mathcal{N}(0,1)$.
Define
\[
X_{n+1} = \rho X_n + \sqrt{1-\rho^2}\,\varepsilon_{n+1}.
\]
Then $K(x,\cdot)=\mathcal{N}(\rho x,\,1-\rho^2)$.
Iterating the recursion shows that for every deterministic starting point \(X_0=x_0\),
\[
  X_n \sim \mathcal{N}\!\left(\rho^n x_0,\,1-\rho^{2n}\right).
\]
In particular, if \(X_0=0\), then \(X_n\sim \mathcal{N}(0,1-\rho^{2n})\), and if $X_0\sim\mathcal{N}(0,1)$, then $X_n\sim\mathcal{N}(0,1)$ for all $n$, so $\pi=\mathcal{N}(0,1)$ is an invariant distribution.
\end{example}

\begin{example}[Ornstein--Uhlenbeck process]
The Ornstein--Uhlenbeck (OU) process is a continuous-time Markov process with a Gaussian transition kernel.
In a standard parameterization, if $X_s=x$ then
\[
X_{s+t}\mid X_s=x \sim \mathcal{N}\!\left(e^{-t}x,\,1-e^{-2t}\right),
\qquad t\ge 0.
\]
Equivalently, for a deterministic initial value \(X_0=x_0\),
\[
  X_t \sim \mathcal{N}\!\left(e^{-t}x_0,\,1-e^{-2t}\right).
\]
Thus from \(X_0=0\) the transient marginal is \(\mathcal{N}(0,1-e^{-2t})\), which relaxes to stationarity as \(t\to\infty\).
The corresponding stationary distribution is $\mathcal{N}(0,1)$.
\end{example}

\subsubsection{Invariant distributions and Monte Carlo averages}
\label{subsec:invariant-stationary}

The central design goal in MCMC is to build a kernel $K$ whose invariant distribution is the target distribution $\pi$.
In terms of the operator $\mu\mapsto \mu K$ defined by
\[
(\mu K)(A)\coloneqq \int_{\mathcal{X}} K(x,A)\,\mu(dx),
\]
invariance simply means that $\pi$ is a fixed point: $\pi K=\pi$.

\begin{definition}[Invariant distribution]
\label{def:invariant-distribution}
A probability measure $\pi$ on $\mathcal{X}$ is \emph{invariant} (or \emph{stationary}) for $K$ if for every $A\in\mathcal{B}(\mathcal{X})$,
\[
\pi(A) = \int_{\mathcal{X}} K(x,A)\,\pi(dx).
\]
Equivalently, if $X\sim \pi$ and $Y\mid X=x \sim K(x,\cdot)$, then $Y\sim \pi$.
\end{definition}

\begin{proposition}[Starting from $\pi$ gives stationarity]
\label{prop:stationarity}
If $(X_n)_{n\ge 0}$ is a Markov chain with kernel $K$ and invariant distribution $\pi$, and $X_0\sim\pi$, then $X_n\sim\pi$ for all $n\ge 0$.
\end{proposition}
\begin{proof}
The claim holds for $n=0$ by assumption.
Assume $X_n\sim\pi$.
Then for any measurable $A$,
\[
\mathbb{P}(X_{n+1}\in A)
= \E\bigl[\mathbb{P}(X_{n+1}\in A\mid X_n)\bigr]
= \E\bigl[K(X_n,A)\bigr]
= \int_{\mathcal{X}} K(x,A)\,\pi(dx)
= \pi(A),
\]
where the last equality is invariance.
\end{proof}

Stationarity means that all one-dimensional marginals are $\pi$, but the samples are generally \emph{not} independent.
This dependence is central to how we analyze the efficiency of MCMC estimators.

\subsubsection{Reversibility and detailed balance}
\label{subsubsec:detailed-balance}
Many MCMC kernels are constructed by enforcing a stronger symmetry property than invariance.

\begin{definition}[Detailed balance]
A kernel $K$ is \emph{reversible} with respect to $\pi$ if for all measurable sets $A,B\in\mathcal{B}(\mathcal{X})$,
\[
\int_A K(x,B)\,\pi(dx) = \int_B K(x,A)\,\pi(dx).
\]
Equivalently, the joint distribution of $(X_0,X_1)$ is symmetric when $(X_n)$ is stationary: if $X_0\sim\pi$ and $X_1\mid X_0\sim K(X_0,\cdot)$, then $(X_0,X_1)$ has the same law as $(X_1,X_0)$.
\end{definition}

Detailed balance implies invariance by taking $A=\mathcal{X}$.
Reversibility is not necessary for MCMC, but it is a convenient sufficient condition and will be the main tool behind Metropolis--Hastings kernels.

\subsubsection{Autocorrelation and effective sample size}
\label{subsubsec:ess}
Let $(X_n)_{n\ge 0}$ be stationary with $X_n\sim\pi$ and let $g:\mathcal{X}\to\R$ satisfy $\E_\pi[g(X)^2]<\infty$.
We estimate $\mu=\E_\pi[g(X)]$ by the empirical average
\[
\hat\mu_n \coloneqq \frac{1}{n}\sum_{i=1}^n g(X_i).
\]
Then $\E[\hat\mu_n]=\mu$, but $\operatorname{Var}(\hat\mu_n)$ is typically larger than the i.i.d.\ variance $\operatorname{Var}_\pi(g(X))/n$ because of serial correlation.

Define the autocovariance $\gamma_k=\operatorname{Cov}_\pi(g(X_0),g(X_k))$.
By expanding the variance of a sum and using stationarity,
\begin{equation}
  \operatorname{Var}(\hat\mu_n)
  = \frac{1}{n^2}\sum_{i=1}^n\sum_{j=1}^n \operatorname{Cov}\!\bigl(g(X_i),g(X_j)\bigr)
  = \frac{1}{n^2}\left(n\gamma_0 + 2\sum_{k=1}^{n-1} (n-k)\gamma_k\right).
  \label{eq:mcmc-variance}
\end{equation}
In the i.i.d.\ case, the off-diagonal covariances vanish and $\operatorname{Var}(\hat\mu_n)=\gamma_0/n$.
This suggests defining the \emph{effective sample size} by matching variances:
\[
N_{\mathrm{eff}}(n,g) \coloneqq \frac{\operatorname{Var}_\pi(g(X_0))}{\operatorname{Var}(\hat\mu_n)}
= \frac{\gamma_0}{\operatorname{Var}(\hat\mu_n)}.
\]
When the samples are i.i.d., $N_{\mathrm{eff}}(n,g)=n$.
When there is positive serial correlation, $N_{\mathrm{eff}}(n,g)$ can be much smaller than $n$.

Writing $\rho_k=\gamma_k/\gamma_0$ (when $\gamma_0>0$) gives
\[
\operatorname{Var}(\hat\mu_n)
= \frac{\gamma_0}{n}\left(1 + 2\sum_{k=1}^{n-1}\left(1-\frac{k}{n}\right)\rho_k\right).
\]
If the series $\sum_{k\ge 1}\rho_k$ is summable, the coefficient in parentheses stabilizes as $n$ grows.
This motivates the \emph{integrated autocorrelation time}
\[
\tau_{\mathrm{int}} \coloneqq 1+2\sum_{k=1}^{\infty}\rho_k,
\]
and the corresponding \emph{effective sample size}
\[
N_{\mathrm{eff}} \approx \frac{n}{\tau_{\mathrm{int}}}.
\]
The interpretation is that $n$ correlated MCMC draws provide roughly the same estimation accuracy as $N_{\mathrm{eff}}$ i.i.d.\ samples from $\pi$.
Both $\tau_{\mathrm{int}}$ and $N_{\mathrm{eff}}$ depend on the function $g$; a chain can mix quickly for some summaries of the state while remaining highly correlated for others.
This is a different notion from the weight-based ESS in importance sampling and SMC.
Here ESS measures the loss of information caused by serial correlation in one chain, whereas there it measures how uneven a set of normalized weights has become.
This autocorrelation-based ESS formula is standard in MCMC diagnostics; see \citet[\S2.2, p.~11]{betancourt_2017_hmc_conceptual}.

\begin{example}[Integrated autocorrelation time for the AR(1) chain]
Consider the stationary AR(1) chain from \Cref{ex:ar1-kernel} and take $g(x)=x$.
Then $\rho_k=\rho^k$, so
\[
\tau_{\mathrm{int}} = 1+2\sum_{k=1}^{\infty}\rho^k = 1+\frac{2\rho}{1-\rho} = \frac{1+\rho}{1-\rho}.
\]
Thus $N_{\mathrm{eff}}\approx n(1-\rho)/(1+\rho)$.
When $\rho$ is close to $1$, successive states are extremely similar and the effective sample size can be much smaller than $n$.
\end{example}

\subsubsection{Ergodicity, mixing time, and burn-in}
\label{subsec:ergodicity}

In practice we rarely start the chain exactly from \(\pi\).
Early iterates still remember the initialization, so averages formed too soon can be biased by where the chain started rather than by the target law itself.
Burn-in is the attempt to discard enough of those early steps that this memory has mostly washed out.
To make that idea precise, we ask that \(K^m(x,\cdot)\) approach \(\pi\) as \(m\to\infty\) for a wide range of starting states \(x\).
Invariance alone does not guarantee this: a chain can have \(\pi\) invariant but fail to reach it from some initial states, or oscillate without settling down.

\begin{definition}[Total variation distance]
\label{def:tv-distance}
For probability measures $\mu,\nu$ on $(\mathcal{X},\mathcal{B}(\mathcal{X}))$, their total variation distance is
\[
\norm{\mu-\nu}_{\mathrm{TV}} \coloneqq \sup_{A\in\mathcal{B}(\mathcal{X})}\bigl|\mu(A)-\nu(A)\bigr|.
\]
\end{definition}

Total variation is a worst-case discrepancy over events: it controls how much any probability statement can change when we replace $\mu$ by $\nu$.
If $\mu$ and $\nu$ have densities $p$ and $q$, then
\[
\norm{\mu-\nu}_{\mathrm{TV}} = \frac12\int_{\mathcal{X}} |p(x)-q(x)|\,dx.
\]

A standard way to summarize convergence speed from a particular starting state $x$ is the $\varepsilon$-mixing time
\[
t_{\mathrm{mix}}(\varepsilon;x) \coloneqq \min\left\{m\ge 0:\norm{K^m(x,\cdot)-\pi}_{\mathrm{TV}}\le \varepsilon\right\}.
\]
On unbounded state spaces, the more stringent ``worst-case over $x$'' notion can be infinite; what matters in practice is convergence from the initializations we actually use.
Under geometric ergodicity (\Cref{def:geometric-ergodicity}), one has
$\norm{K^m(x,\cdot)-\pi}_{\mathrm{TV}}\le M(x)\rho^m$, so it is enough to take $m$ on the order of $\log(M(x)/\varepsilon)/\log(1/\rho)$.
The burn-in length $B$ in \Cref{alg:mcmc-burnin} is an attempt to choose $B\gtrsim t_{\mathrm{mix}}(\varepsilon;x^{(0)})$ for some small $\varepsilon$, but in most realistic problems we cannot compute $t_{\mathrm{mix}}$ and instead rely on diagnostic checks.
This separation between early bias reduction and later variance reduction is the same qualitative warm-up picture emphasized in \citet[\S2.2, pp.~11--12]{betancourt_2017_hmc_conceptual}.

\begin{definition}[Geometric ergodicity]
\label{def:geometric-ergodicity}
A Markov chain with transition kernel $K$ and invariant distribution $\pi$ is \emph{geometrically ergodic} if there exist $\rho\in(0,1)$ and a function $M:\mathcal{X}\to[0,\infty)$ such that for all $m\ge 0$,
\[
\norm{K^m(x,\cdot)-\pi}_{\mathrm{TV}} \le M(x)\rho^m.
\]
\end{definition}

Geometric ergodicity is a strong form of convergence: it says the distribution contracts at a linear (geometric) rate in total variation.
The Gaussian AR(1) chain in \Cref{ex:ar1-kernel} and the OU process are clean examples of this behavior: their linear drift contracts toward the origin while the noise prevents collapse.

\begin{algorithm}[t]
  \caption{Generic MCMC estimation with burn-in}
  \label{alg:mcmc-burnin}
  \begin{algorithmic}[1]
    \Require Transition kernel $K$ with invariant distribution $\pi$, initial state $x^{(0)}$, burn-in $B$, total steps $T$, test function $g$
    \For{$t=0,1,\dots,T-1$}
      \State Draw $x^{(t+1)} \sim K(x^{(t)},\cdot)$
    \EndFor
    \State Discard $x^{(0)},\dots,x^{(B)}$ and keep $x^{(B+1)},\dots,x^{(T)}$
    \State \Return $\hat\mu = (T-B)^{-1}\sum_{t=B+1}^T g(x^{(t)})$
  \end{algorithmic}
\end{algorithm}

Discarding burn-in reduces bias from initialization, but it does not make the remaining samples independent.
The residual serial dependence is what drives the variance inflation in \Cref{eq:mcmc-variance} and is summarized by $\tau_{\mathrm{int}}$ (or $N_{\mathrm{eff}}$).

When comparing MCMC algorithms, two recurring questions are:
(i) how long it takes for the chain to approach stationarity (mixing time / burn-in), and
(ii) how much correlation remains after convergence (which controls $N_{\mathrm{eff}}$ via \Cref{eq:mcmc-variance}).
We return to practical diagnostics for both questions in \Cref{subsec:mcmc-diagnostics}.
